\clearpage\thispagestyle{empty}
\pdfbookmark[0]{Quellen zu den Neuen Pallas-Elementen}{pallaselements1810oct-sources}
\begin{center}\textbf{Quellen zu den Neuen Pallas-Elementen}\end{center}
\noindent Diese KI-Übertragung ins heutige Deutsch folgt dem Beitrag \textit{Neue Pallaselemente} in \href{https://gdz.sub.uni-goettingen.de/id/PPN236007467}{Gauß’ \textit{Werken}, Bd.~VI} (Göttingen, 1874), S.~320. Beitrag XLV der Zeitschrift \href{https://archive.org/download/bub_gb_YCI4AAAAMAAJ/page/n4.jpg}{\textit{Monatliche Correspondenz}, Bd.~XXII} vom Oktober 1810 nennt Gauß auf S.~\href{https://archive.org/download/bub_gb_YCI4AAAAMAAJ/page/n403.jpg}{400} und druckt seine Ephemeride auf S.~\href{https://archive.org/download/bub_gb_YCI4AAAAMAAJ/page/n404.jpg}{401}, \href{https://archive.org/download/bub_gb_YCI4AAAAMAAJ/page/n405.jpg}{402} und \href{https://archive.org/download/bub_gb_YCI4AAAAMAAJ/page/n406.jpg}{403}. Das digitalisierte Exemplar der University of Michigan ist im Internet Archive zugänglich.
\par\smallskip\noindent Die Werkausgabe bringt den Bericht und die Bahnelemente. Es folgt eine eckige Klammer, die die Ephemeride beschreibt. Im Bericht steht das redaktionelle Wir. Gauß wird in der dritten Person als Urheber der Berechnung genannt. Die Zeile des aufsteigenden Knotens, die kommaförmige Markierung in der Zeile der Neigung und die Winkelzeichen folgen der Werkausgabe. KI-Übertragung und KI-Satz.
\clearpage
